% 366_Feldstruktur_FFGFT_En_ch.tex
% Johann Pascher, 19 September 2026
% Fields in FFGFT: Energy, Flux and Geometry

\providecommand{\BM}{\textbf{[B]}}
\providecommand{\KM}{\textbf{[K]}}
\providecommand{\SM}{\textbf{[S]}}
\providecommand{\SetzM}{\textbf{[POSIT]}}
\providecommand{\QM}{\textbf{[Q]}}
\providecommand{\EM}{\textbf{[E]}}
\providecommand{\XM}{\textbf{[X]}}

\chapter*{Fields in FFGFT: Energy, Flux and Geometry}
\addcontentsline{toc}{chapter}{Fields in FFGFT: Energy, Flux and Geometry}

\begin{abstract}
This document develops a unified approach to the field structure of FFGFT,
starting from an unexpected experimental result and ending at the orbital
resonances of the planets. The entry point is deliberately concrete: the
Goubau line experiment --- energy transport along a single wire without
a return conductor --- forces the question of where energy actually flows
to be taken seriously. The answer has been fully present in classical field
theory since Poynting (1884); it is regularly obscured in introductory
courses, however, by the useful but simplified picture of current and
return path.

The goal is not to claim new physics. The goal is to make visible that
antenna engineering, atomic physics, crystallography and celestial mechanics
all realise structurally the same principle: the geometry of the carrier
determines the permitted coupling angles. Coupling between two field systems
is only possible where their mode structures are compatible. Incompatible
modes do not couple --- regardless of the field strength.

In FFGFT this connection is not an analogy but a consequence of the
fundamental equation $\tilde{T}\cdot m = 1$ on the compact carrier
$T^4/\mathbb{Z}_3$. The winding numbers $(n_\theta, n_\varphi)$ on this
carrier are the topological invariants behind the selection rules of atomic
physics, the Bragg conditions of crystallography and the integer orbital
resonances of celestial mechanics. In FFGFT units ($c=\hbar=\alpha=1$):
$I=m$, $V=E$, $P=E\cdot m$ --- exactly, not approximately.

Antenna engineering serves as the guide because it is the most mature
practical application of electromagnetic field theory: the coupling
conditions are not abstract there but codified in standards, measurement
protocols and commercial simulators. What is known there as the Friis
equation, impedance matching and array factor appears in atomic physics
as a selection rule, in the crystal as the Bragg condition and in the
solar system as the Laplace resonance --- the same structure, different carriers.
\end{abstract}

% ============================================================
\section*{What this document explains --- and what it does not}
\addcontentsline{toc}{section}{What this document explains --- and what it does not}

Four questions are addressed, developed in this order:

\textbf{1. Where does energy flow?} Not in the conductor, but in the field
around it. The Goubau experiment forces this answer experimentally; the
Poynting vector delivers it mathematically exactly. Current and voltage
are projections, not primary quantities (\S\,1--3).

\textbf{2. What does an antenna do to the field?} It actively reshapes
it --- in transmission by imposing boundary conditions, in reception by
acting as a selective mode filter. Passive in a field it scatters; in
near-field proximity two antennas mutually deform each other. This is
not superposition but mutual deformation of mode structures (\S\,4--9).

\textbf{3. What is a particle in this picture?} A stable field node ---
a localised mode concentration on the carrier. The near field of the
particle is the Coulomb field; the far field arises only under
acceleration or a state transition. The selection rules of atomic
physics are the mode-compatibility conditions of this node (\S\,10).

\textbf{4. Does the principle hold on all scales?} Yes --- from the
antenna through the atom and the crystal to the planetary orbits. The
coupling condition is always the same: integer phase compatibility of
the mode spectra. The scale changes; the structure of the condition
does not (\S\,11).

What this document does \emph{not} explain: why the carrier is
$T^4/\mathbb{Z}_3$ (that is \SetzM\ in Doc.~365); why $\xi = 4/30\,000$
(that follows from Galois group orders, Doc.~338/365); how concrete
masses follow from mode compatibility (Doc.~314, 317--323). Those
documents are the sources cited here. This document is the approach
from the concrete to the structural.

\vspace{0.5em}
\noindent
\textit{Reading note:} Readers familiar with antenna engineering may skim
\S\,1--3 and enter at \S\,4. Those who know the FFGFT background will find
the connection to mode structure in \S\,5, \S\,10 and \S\,11.

% ============================================================
\section*{List of symbols}
\addcontentsline{toc}{section}{List of symbols}

\medskip
\textbf{Epistemic markers:}
\begin{tabular}{@{}ll@{}}
\textbf{[B]} & algebraically proven (Python script, exact arithmetic) \\
\textbf{[K]} & numerically confirmed \\
\textbf{[S]} & structurally plausible, identification open \\
\textbf{[POSIT]} & motivated input, not derived \\
\textbf{[E]} & established external mathematics or physics \\
\textbf{[X]} & checked and negative \\
\end{tabular}

\medskip
\begin{longtable}{@{}lp{11cm}@{}}
\toprule
\textbf{Symbol} & \textbf{Meaning} \\
\midrule
\multicolumn{2}{l}{\textit{Fields and energy}} \\
$F^{\mu\nu}$ & electromagnetic field strength tensor \\
$\mathbf{E}, \mathbf{B}$ & electric and magnetic field \\
$\mathbf{H}$ & magnetic field strength ($\mathbf{H} = \mathbf{B}/\mu_0$ in vacuum) \\
$\mathbf{S} = \mathbf{E}\times\mathbf{H}$ & Poynting vector (energy flux density, W/m$^2$) \\
$j^\nu$ & electromagnetic four-current \\
$u$ & electromagnetic energy density \\
$\varepsilon_0,\,\mu_0$ & permittivity of free space, permeability of free space \\
$\lambda$ & wavelength ($\lambda = c/f$) \\
$\omega = 2\pi f$ & angular frequency \\
$k = 2\pi/\lambda$ & wave number \\
$\theta$ & elevation angle (radiation pattern) or glancing angle (Bragg) \\
$\psi$ & polarisation mismatch angle between transmit and receive antenna \\
\midrule
\multicolumn{2}{l}{\textit{FFGFT quantities}} \\
$\tilde{T}$ & intrinsic time scale; $\tilde{T} \cdot m = 1$ \\
$\xi = 4/30\,000$ & dimensionless scaling parameter of FFGFT \\
$T^4/\mathbb{Z}_3$ & compact orbifold carrier of FFGFT \\
$(n_\theta, n_\varphi)$ & winding numbers of a field mode on $T^4$; topological invariants \\
\midrule
\multicolumn{2}{l}{\textit{Antenna engineering}} \\
$R_\mathrm{rad}$ & radiation resistance (resistance the antenna presents for radiated power) \\
$G$ & antenna gain (relative to isotropic radiator) \\
$D$ & directivity (gain without losses) \\
$A_\mathrm{eff}$ & effective aperture; $A_\mathrm{eff} = (\lambda^2/4\pi)\,G$ \\
$\eta_A$ & aperture efficiency ($0 < \eta_A \leq 1$) \\
$G_T, G_R$ & gain of transmit and receive antenna (Friis equation) \\
$\Delta\theta_{3\,\mathrm{dB}}$ & half-power beamwidth of main lobe \\
$Z_\mathrm{ant}$ & input impedance of the antenna \\
$X_\mathrm{ant}$ & reactance (imaginary part of $Z_\mathrm{ant}$) \\
$\phi_n$ & feed phase of the $n$-th element in a phased array \\
$AF(\theta)$ & array factor (interference pattern of the phased array) \\
$\varepsilon_r$ & relative permittivity (dielectric constant) \\
$\varepsilon_\mathrm{eff}$ & effective permittivity of a microstrip line \\
$\Delta L$ & fringe correction of the patch antenna length \\
$k_\mathrm{SPP}$ & wave number of a surface plasmon polariton \\
$\varepsilon_m, \varepsilon_d$ & (complex) permittivity of metal and dielectric (SPP) \\
\midrule
\multicolumn{2}{l}{\textit{Atomic physics and scales}} \\
$\hbar$ & reduced Planck constant \\
$\alpha \approx 1/137$ & fine-structure constant; in FFGFT units $\alpha = 1$ \\
$n, l, m$ & principal, orbital angular momentum, and magnetic quantum numbers \\
$E_n$ & energy levels of the hydrogen atom; $E_n = -13.6\,\mathrm{eV}/n^2$ \\
$r_e \approx 2.82\,\mathrm{fm}$ & classical electron radius \\
$\lambda_C = \hbar/m_e c$ & Compton wavelength of the electron (${\approx}386\,\mathrm{fm}$) \\
$a_0 \approx 52\,917\,\mathrm{fm}$ & Bohr radius \\
$\sigma_T$ & Thomson cross-section; $\sigma_T = \tfrac{8\pi}{3}r_e^2$ \\
$S_{AB}$ & overlap integral of two atomic orbitals $\psi_A$, $\psi_B$ \\
$Q$ & normal coordinate of a molecular vibration \\
$\vec{\mu}$ & electric dipole moment of a molecule \\
$d$ & lattice plane spacing in a crystal (Bragg condition) \\
$n_1, n_2, n_3$ & mean motions (orbital frequencies) of the Galilean moons \\
\bottomrule
\end{longtable}

% ============================================================
\section{Starting point: fields or current?}
% ============================================================

Every introductory course on electricity emphasises the principle: current
needs a return path. Electrons leaving one conductor must return on the
other. This picture is entirely correct for technical purposes --- it
describes circuits precisely. It conceals a deeper question, however:
where does the energy actually flow?

The Goubau experiment \QM\ (G.~Goubau, 1950) demonstrates the answer
experimentally. A high-frequency signal source (${\sim}900\,\mathrm{MHz}$)
feeds a single isolated wire via a conical transition (launcher). At the
other end an identical cone (catcher) converts the signal back into a
coaxial line. Between the two cones there is no second conductor, no return
path in the classical sense --- yet the arrangement transfers the signal
with low losses \cite{Goubau1950}.

A physical interruption of the field (for instance a hand placed over the
wire) immediately interrupts the transfer. The energy flows in the field
around the wire, not in the wire itself.

% ============================================================
\section{The Poynting vector as the primary description}
% ============================================================

The energy flux density of electromagnetic fields is the Poynting vector
\cite{Stratton1941,Griffiths1999}:
\[
  \mathbf{S} = \mathbf{E} \times \mathbf{H}.
\]
Its direction (right-hand rule: $\mathbf{E} \to \mathbf{H}$, thumb in the
direction of flux) indicates where energy is transported. The flux density
integrated over the cross-section of a cable gives exactly the electrical
power \EM:
\[
  P = \oint_A \mathbf{S} \cdot d\mathbf{A} = V \cdot I.
\]
This is not an approximation, not a coincidence --- it is an exact identity
\cite{Griffiths1999}. Current and voltage are macroscopic projections of the
field process, not the primary carriers of energy.

\paragraph{Consequence.}
For a coaxial line the $\mathbf{E}$ field lines run radially from the inner
conductor to the outer shield; the $\mathbf{H}$ field lines form concentric
rings. The cross product $\mathbf{E}\times\mathbf{H}$ points axially ---
along the line. Energy flows in the dielectric between inner and outer
conductor, not in the copper.

% ============================================================
\section{The Goubau line: fields without a return conductor}
% ============================================================

The Goubau line (G-line) has no outer conductor. The field configuration
differs from the coaxial line only gradually when measured close enough to
the wire \EM.

\subsection*{Field configuration}

A voltage pulse propagating on the wire creates alternating positive and
negative potential regions along the conductor. The electric field lines
run between these regions --- not from wire to return conductor, but from
one section of the wire to the next. The dielectric (the red PVC insulation)
slows the propagation of the fields and bends the field lines towards the
wire axis, just as glass bends light rays \cite{Goubau1950,Harms1907}.
This mechanism keeps the fields close to the wire --- it is not cosmetic
but physically necessary for practical operation.

\subsection*{The cone transition as impedance matching}

The cones (launcher / catcher) match the wave impedance of the coaxial
line ($50\,\Omega$) to the characteristic impedance of the G-line
(typically a few $100\,\Omega$) \EM. Inside the cone the field lines
expand from the coaxial pattern to the G-line pattern. At the opening of
the cone the field lines snap and attach to the wire. The return path for
the current in the coaxial cable is provided by the displacement currents
inside the cone --- the complex fields in the cone supply exactly the right
forces to drive electrons in the outer surface.

\subsection*{Losses}

Because current flows in the wire, $I^2 R$ losses occur --- exactly as in
a conventional two-wire line. This shows: the G-line is not a pure
waveguide problem. The energy flux has a field component (Poynting) and a
conduction component (Ohm); both describe the same physics from different
perspectives.

% ============================================================
\section{FFGFT units: current is mass, voltage is energy}
% ============================================================

In FFGFT units ($c = \hbar = \alpha = 1$, Doc.~011; Doc.~365 Layer~2)
electromagnetic quantities simplify to geometric expressions in mass and
time \KM:
\begin{align*}
  e &= 1 \quad \text{(elementary charge dimensionless)},\\
  V &= E/e = E \quad \text{(voltage = energy)},\\
  I &= e/T = 1/T = m \quad \text{(current = mass, via }\tilde{T}\cdot m=1\text{)}.
\end{align*}
The power equation then reads:
\[
  P = V \cdot I = E \cdot m,
\]
a relation that follows directly from the fundamental equation
$\tilde{T}\cdot m=1$ (Doc.~365 Layer~1) \KM.

\paragraph{Maxwell equations in FFGFT units.}
In FFGFT units (Doc.~365 \S\,2) the Maxwell equations take the form:
\[
  \partial_\mu F^{\mu\nu} = j^\nu, \qquad \partial_{[\mu}F_{\nu\rho]}=0.
\]
No prefactors, no coupling constants other than $j^\nu$ as a geometric
object. The coupling is absorbed in the unit choice, not eliminated; the
physical content is identical to the SI system.

% ============================================================
\section{Field lines as geometric objects on $T^4/\mathbb{Z}_3$}
% ============================================================

In FFGFT the physical carrier is the compact orbifold $T^4/\mathbb{Z}_3$
(Doc.~365 Layer~1) \SetzM. Fields are sections of vector bundles over this
carrier; their periodicity conditions enforce discrete mode spectra.

\subsection*{Field lines and winding numbers}

On a torus $T^4$ non-trivial field configurations have topological
invariants --- winding numbers $(n_\theta, n_\varphi)$ that record how many
times a field line traverses the torus in each direction (Doc.~314;
Doc.~247) \BM. These winding numbers classify the modes and determine ---
via $\tilde{T}\cdot m=1$ --- the mass scale of the respective sector. A
field line in this description is not an auxiliary construct but a
topological object with physical relevance.

\subsection*{Field lines in the G-line context}

The behaviour of field lines in the Goubau line --- expansion in the cone,
snapping at the opening, attachment to the wire --- is from this perspective
a macroscopic example of the energy-flux minimisation principle. At every
position and time the flux chooses the path that keeps the Poynting vector
consistent with the carrier. This is not mysterious behaviour but the
continuum version of topological optimisation on $T^4/\mathbb{Z}_3$.

\subsection*{Field shape and carrier: no one-to-one correspondence}

The fundamental equation $\tilde{T}\cdot m=1$ says nothing about the
\emph{shape} of a field in space. It specifies which mass scales belong
to which time scales --- not how the field looks geometrically. The compact
carrier $T^4/\mathbb{Z}_3$ is the topology that prescribes the permitted
boundary conditions; the actual field configurations on and around this
carrier are different from it and variously deformed depending on the
situation.

A comparison clarifies this: the length of a guitar string determines the
permitted frequency spectrum --- but the sound wave in the surrounding
space is not string-shaped. Equally, $T^4/\mathbb{Z}_3$ enforces the mode
spectrum; what appears as a field in the three-dimensional observation space
is a projection of it, depending on the concrete boundary conditions,
the ambient geometry and the mode occupation. Close to the conductor the
G-line field is cylindrical, further away elliptically deformed, at
inhomogeneities asymmetric --- none of this is torus-shaped. The torus is
the internal geometric carrier, not the shape of the field.

\subsection*{Interaction: carrier determines permitted coupling angles}

When two fields from different carriers overlap in space, their interaction
is not simple superposition. Each field carries its own mode structure,
fixed by the winding numbers $(n_\theta, n_\varphi)$ of its carrier.
Coupling is only possible where the phase conditions of both mode spectra
are compatible --- where the winding numbers fit together as integers. The
geometry of the carrier thus determines the \emph{angles} of interaction:
under which phase relations energy transfer occurs and under which it does
not \SM.

This is worked out quantitatively in the corpus: the Galois lattice
condition from Doc.~328 (coupling regimes) and Doc.~343~\S\,G$'$ specifies
that coupling is only effective at frequency ratios in the Galois lattice
$\{2,3,5,7,11,13\}$; the resolution floor $100\xi$ is the sharpest bound
below which two modes are no longer distinguishable as separate. The
apparently free shape of a field in space is the result of these
compatibility conditions, projected onto the observation space \SM.

% ============================================================
\section{Antenna engineering as mature field theory}
% ============================================================

Antenna engineering is the most mature practical application of
electromagnetic field theory. Decades of measurements, standards and
calculation procedures are available; the coupling conditions are not
abstract but anchored in handbooks, measurement protocols and commercial
simulators. It is precisely for this reason the most suitable introduction
to the general principle: the geometry of the carrier determines the
permitted coupling angles.

\subsection*{Fundamental quantities and their field interpretation}

Every antenna is a transducer between a guided wave (coaxial cable,
waveguide) and a free wave in space. The decisive quantities are \EM:

\begin{itemize}[nosep]
\item \textbf{Radiation resistance} $R_\mathrm{rad}$: The resistance the
  antenna presents to the guided wave because it radiates energy into space.
  For the half-wave dipole: $R_\mathrm{rad} \approx 73.1\,\Omega$.
  In FFGFT units ($I=m$, $V=E$): $R_\mathrm{rad} = P/I^2 = E/m$.
  The radiation resistance is thus an energy-per-mass relation --- a direct
  echo of $\tilde{T}\cdot m=1$ \SM.

\item \textbf{Gain} $G$ and \textbf{directivity} $D$: Gain gives the factor
  by which the radiation intensity in the main direction exceeds that of an
  isotropic radiator. It is a pure field-geometry problem: how do the phase
  conditions of the carrier (antenna length, element spacing, feed phase)
  concentrate the Poynting vector in a preferred direction? For the half-wave
  dipole: $D = 1.64$ (2.15\,dBi).

\item \textbf{Effective aperture} $A_\mathrm{eff}$: The area from which a
  receiving antenna draws energy. Relation to gain \EM:
  \[
    A_\mathrm{eff} = \frac{\lambda^2}{4\pi}\,G.
  \]
  This formula is pure field geometry: wavelength and gain determine the
  effective area, not the physical size of the antenna.
\end{itemize}

\subsection*{Friis transmission equation: coupling of two carriers}

The Friis equation describes energy transfer between two antennas in free
space \cite{Friis1946,Balanis2016}:
\[
  \frac{P_R}{P_T} = G_T\,G_R\left(\frac{\lambda}{4\pi r}\right)^2
  \cos^2\psi,
\]
where $P_T$ is the transmitted, $P_R$ the received power, $G_T, G_R$ the
gains of transmit and receive antennas, $r$ the distance and $\psi$ the
polarisation mismatch angle.

Every factor of this equation corresponds to a compatibility condition
between the carriers:

\begin{center}
\begin{tabular}{@{}lll@{}}
\toprule
\textbf{Factor} & \textbf{Physical meaning} & \textbf{FFGFT correspondence} \\
\midrule
$G_T$ & focussing of transmit carrier & mode geometry of source field \\
$G_R$ & focussing of receive carrier & mode geometry of receive field \\
$(\lambda/4\pi r)^2$ & free-space attenuation & geometric propagation factor \\
$\cos^2\psi$ & polarisation overlap & winding-direction compatibility \\
\bottomrule
\end{tabular}
\end{center}

The $\cos^2\psi$ factor is the most direct expression of the principle: at
orthogonal polarisation ($\psi=90°$) transfer is zero --- not because of
absorption or loss but because the mode geometries are incompatible. No
transmitter of arbitrary power overcomes this. This is identical to the
statement about winding numbers: if they do not match, no energy transfer
occurs \SM.

\subsection*{Impedance matching as a phase condition}

For an antenna to radiate at maximum efficiency its input impedance must
be the complex conjugate of the source impedance (conjugate matching) \EM:
\[
  Z_\mathrm{ant} = R_\mathrm{rad} + jX_\mathrm{ant}, \qquad
  Z_\mathrm{src} = R_\mathrm{src} - jX_\mathrm{ant}.
\]
The imaginary part --- the reactance --- measures stored field energy near
the antenna (near field, reactive power). If it is not compensated, energy
oscillates between source and antenna near field without being radiated.
This is the same phenomenon as the reactive power on the G-line (\S\,3):
energy pulses in the near field without net transport.

For the G-line the cone transition was the matching structure
($50\,\Omega$ to ${\sim}300\,\Omega$); for antennas this is handled by
$\lambda/4$ transformers, baluns or matching networks. The mechanism is
structurally identical: transition between two mode spectra through
geometrically enforced phase condition.

\subsection*{Antenna types as mode geometries}

Different antenna types realise different mode geometries. Each type is
an independent carrier with a characteristic mode spectrum; the radiation
pattern in the far field is the projection of this spectrum onto the
observation space.

% -----------------------------------------------------------
\paragraph{1. Short dipole ($\ell \ll \lambda$).}

The short dipole is the simplest limiting case: a straight conductor whose
length is far below the wavelength. The current distribution is
approximately triangular; the far field is analytically exact \EM:
\[
  E_\theta \propto \frac{I_0\,\ell}{\lambda r}\,\sin\theta,
  \qquad
  F(\theta) = \sin^2\theta.
\]
The radiation pattern $F(\theta) = \sin^2\theta$ is a torus shape around
the dipole axis: maximum radiation perpendicular to the axis
($\theta = 90°$), zero along the axis ($\theta = 0°, 180°$). Directivity
$D = 1.5$ (1.76\,dBi). Radiation resistance:
\[
  R_\mathrm{rad} = 80\pi^2\left(\frac{\ell}{\lambda}\right)^2 \approx 790\,\Omega
  \cdot \left(\frac{\ell}{\lambda}\right)^2.
\]
For $\ell = 0.1\lambda$: $R_\mathrm{rad} \approx 7.9\,\Omega$ --- very low,
hence severe mismatch to $50\,\Omega$ systems. The mode structure is simple:
a single electric dipole mode whose winding geometry enforces the torus
shape of the radiation pattern.

% -----------------------------------------------------------
\paragraph{2. Half-wave dipole ($\ell = \lambda/2$).}

The half-wave dipole is the most important reference radiator in antenna
engineering. The current distribution is sinusoidal; the radiation pattern
is slightly narrower than the short dipole \EM:
\[
  F(\theta) = \frac{\cos^2\!\left(\frac{\pi}{2}\cos\theta\right)}{\sin^2\theta}.
\]
Directivity $D = 1.64$ (2.15\,dBi). Radiation resistance
$R_\mathrm{rad} = 73.1\,\Omega$ --- nearly ideal for $50\,\Omega$ or
$75\,\Omega$ coaxial cable (mismatch $< 2.5:1$ without matching network).
The half-wave resonance condition ($\ell = \lambda/2$) enforces exactly
the phase relationship at which radiation resistance and feed resistance
enter a real ratio: the imaginary part of the input impedance vanishes
at resonance. This is an integer mode-compatibility condition --- the
length must be a half-integer multiple of the wavelength.

% -----------------------------------------------------------
\paragraph{3. Monopole ($\ell = \lambda/4$) over ground plane.}

A quarter-wave monopole over an infinitely extended conducting plane is
equivalent to the half-wave dipole by image theory: the ground plane
supplies the mirror image. The radiation pattern is the lower hemisphere
of the dipole torus, reflected upward; all radiated power concentrates in
the upper half-space \EM:
\[
  D_\mathrm{mono} = 2\,D_\mathrm{dipole} = 3.28 \quad (5.15\,\mathrm{dBi}).
\]
Radiation resistance $R_\mathrm{rad} = 36.6\,\Omega$ (half that of the
dipole, since only half the space). Practical implementations: car antenna,
mobile mast \cite{Balanis2016}. The gain increase over the dipole arises not
from more input power but from the geometric mirroring condition ---
doubling the solid angle contributing to radiation is a pure mode-geometry
question.

% -----------------------------------------------------------
\paragraph{4. Yagi--Uda antenna.}

The Yagi--Uda antenna consists of a driven dipole (radiator), a slightly
longer passive reflector and a series of shorter passive directors
\cite{Yagi1928,Balanis2016}. All parasitic elements are excited by induction;
their length determines the phase of the induced current:
\begin{itemize}[nosep]
  \item Reflector ($\ell \approx 0.5\lambda + 5\,\%$): capacitive,
    re-radiates with phase lag $\Rightarrow$ cancellation to the rear.
  \item Directors ($\ell \approx 0.5\lambda - 5\,\%$): inductive,
    re-radiate with phase lead $\Rightarrow$ constructive interference
    forward.
\end{itemize}
Gain as a function of the number of directors $n$ (rule of thumb) \EM:
\[
  G \approx 10\log_{10}(0.8\,n + 1)\;\mathrm{dBd}
  \quad (n = \text{number of directors}).
\]
Example values: 3 directors $\Rightarrow G \approx 7\,\mathrm{dBd}$,
10 directors $\Rightarrow G \approx 10\,\mathrm{dBd}$.
The half-power beamwidth (in the $E$ plane) is approximately:
\[
  \Delta\theta_{3\,\mathrm{dB}} \approx \frac{52°}{\sqrt{G_\mathrm{lin}}},
\]
where $G_\mathrm{lin}$ is the linear gain factor. The radiation pattern
narrows with increasing element count --- the carrier becomes more selective
in its coupling angles.

% -----------------------------------------------------------
\paragraph{5. Horn antenna.}

The horn antenna is the transition from waveguide to free space. The
waveguide supports a discrete mode structure (TE$_{10}$ fundamental mode
for rectangular waveguide); the horn adiabatically expands this mode into
a plane wave in the far field \cite{Balanis2016,Pozar2012}. The half-power
beamwidth depends directly on the aperture dimensions:
\[
  \Delta\theta_{E} \approx 56°\,\frac{\lambda}{a}, \qquad
  \Delta\theta_{H} \approx 67°\,\frac{\lambda}{b},
\]
where $a$ is the $E$-plane aperture and $b$ the $H$-plane aperture.
Gain of a pyramidal horn with aperture $A = a \times b$:
\[
  G = \eta_A\,\frac{4\pi A}{\lambda^2}, \qquad \eta_A \approx 0.5\text{--}0.7.
\]
The aperture efficiency $\eta_A < 1$ arises from the non-uniform phase
distribution in the aperture (spherical rather than plane wave at the
horn mouth). This is directly the mode incompatibility between waveguide
near field and far-field plane wave: the better the phase matching (flatter
aperture), the higher $\eta_A$.

% -----------------------------------------------------------
\paragraph{6. Parabolic antenna.}

A parabolic antenna of diameter $D$ and focal length $f$ focuses incident
microwaves at the focus. In transmission a primary radiator at the focus
converts a spherical wave into a plane wavefront at the aperture plane.
The far field is the Fourier transform of this aperture distribution \EM.
For a circular, uniformly illuminated aperture:
\[
  F(\theta) = \left[\frac{2\,J_1(\pi D\sin\theta/\lambda)}
                        {\pi D\sin\theta/\lambda}\right]^2,
\]
where $J_1$ is the first-order Bessel function \cite{Balanis2016}. Gain
and half-power beamwidth:
\[
  G = \eta_A\,\frac{\pi^2 D^2}{\lambda^2}, \qquad
  \Delta\theta_{3\,\mathrm{dB}} \approx 70°\,\frac{\lambda}{D}.
\]
For a 3-m parabolic antenna at 10\,GHz ($\lambda = 3\,\mathrm{cm}$):
$G \approx 10^4$ (40\,dBi), $\Delta\theta \approx 0.7°$. The Bessel
function in the radiation pattern is not an approximation but the exact
consequence of the circular symmetry of the aperture: the mode geometry
(circle) completely determines the coupling pattern in the far field.

% -----------------------------------------------------------
\paragraph{7. Patch antenna (microstrip).}

A rectangular patch antenna of length $L$ and width $W$ over a ground
plane (dielectric $\varepsilon_r$, thickness $h$) resonates when $L$
equals half an effective wavelength \cite{Pozar2012,Balanis2016}:
\[
  L = \frac{\lambda}{2\sqrt{\varepsilon_\mathrm{eff}}} - 2\Delta L,
\]
where $\varepsilon_\mathrm{eff}$ is the effective permittivity and
$\Delta L$ the fringe correction:
\[
  \varepsilon_\mathrm{eff} = \frac{\varepsilon_r+1}{2}
    + \frac{\varepsilon_r-1}{2}\left(1 + \frac{12h}{W}\right)^{-1/2}.
\]
The radiation pattern is broadside in the $H$ plane (no focussing), narrow
in the $E$ plane; typical gain 5--8\,dBi. The permitted modes are
completely determined by geometry and material ($L$, $W$, $h$,
$\varepsilon_r$) --- no free parameters. Coupling to an external wave is
only possible when its polarisation and angle of incidence are compatible
with the resonance mode of the patch. Cross-polarisation isolation
($> 30\,\mathrm{dB}$) shows: orthogonal polarisation practically does
not couple, even when the frequency matches exactly.

% -----------------------------------------------------------
\paragraph{8. Phased array.}

A phased array consists of $N$ identical radiators with individually
controlled feed phase $\phi_n$. The array pattern is the product of the
element radiation pattern $F_\mathrm{elem}(\theta)$ and the array factor
\cite{Mailloux2017,Balanis2016}:
\[
  AF(\theta) = \sum_{n=0}^{N-1} \exp\!\left[j\left(n\,kd\cos\theta + \phi_n\right)\right],
\]
where $k = 2\pi/\lambda$ is the wave number and $d$ the element spacing.
By choosing $\phi_n = -n\,kd\cos\theta_0$ the main beam steers to angle
$\theta_0$ --- electronic beam steering without mechanical movement. Gain
for $N$ elements with spacing $d = \lambda/2$:
\[
  G_\mathrm{array} = N \cdot G_\mathrm{elem}.
\]
The phased array is the purest example of the coupling-angle principle:
the feed phases $\phi_n$ are integer multiples of a basic phase step ---
they determine under which angle constructive interference occurs. Every
other angle combination gives destructive interference. In FFGFT language:
the $\phi_n$ are discrete winding phase values; the array factor selects
the compatible coupling angles \SM.

% -----------------------------------------------------------
\paragraph{Summary table.}

\begin{center}
\begin{tabular}{@{}llll@{}}
\toprule
\textbf{Type} & \textbf{Gain} & \textbf{$\Delta\theta_{3\,\mathrm{dB}}$}
             & \textbf{Mode geometry} \\
\midrule
Short dipole     & 1.76\,dBi  & $90°$         & electric dipole mode \\
Half-wave dipole & 2.15\,dBi & $78°$          & resonance $\ell=\lambda/2$ \\
Monopole ($\lambda/4$) & 5.15\,dBi & $46°$    & image mode over ground \\
Yagi (5 dir.)   & ${\approx}12\,\mathrm{dBd}$ & ${\approx}30°$ & interference array \\
Horn antenna    & 15--25\,dBi & $3°$--$15°$  & TE$_{10}$ aperture mode \\
Parabolic ($D=3\,\mathrm{m}$, 10\,GHz) & 40\,dBi & $0.7°$ & Fourier aperture \\
Patch           & 5--8\,dBi  & $60°$--$120°$ & resonance mode in dielectric \\
Phased array ($N$) & $N \cdot G_\mathrm{elem}$ & $\propto 1/\sqrt{N}$ & phase-steered array \\
\bottomrule
\end{tabular}
\end{center}

In every case: the physical geometry of the antenna determines the permitted
mode spectrum; what appears in the far field as a radiation pattern is the
projection of this spectrum onto the observation space. No antenna type
transfers energy at angles that contradict its mode geometry --- not
approximately but structurally. This is the technically mature, quantitatively
anchored form of the principle formulated in FFGFT as winding-number
compatibility \SM.

\subsection*{Reciprocity: transmitter and receiver as dual descriptions}

The reciprocity principle of antenna engineering states: gain, radiation
pattern and impedance of an antenna are identical in transmit and receive
mode \cite{Balanis2016,Griffiths1999}. This symmetry is not coincidental ---
it reflects the time-reversal symmetry of the Maxwell equations, and with it
the duality of field source and field sink. In FFGFT language: transmitter
and receiver are reciprocal projections of the same field process, just as
$\tilde{T}$ and $m$ are reciprocal representations of the same object \SM.

% ============================================================
\section{Multiple field carriers in space: boundary conditions, coupling and apparent superposition}
% ============================================================

The superposition principle of classical field theory states: in a linear
medium the fields of independent sources add linearly. Each partial field
propagates as if the other were not present. This is an excellent
approximation for widely separated, weakly coupling sources --- and it
works very well in practice. It conceals, however, what physically actually
occurs.

\subsection*{The total field as the consistent solution of all boundary conditions}

Every object in a field --- an antenna, a conductor, a dielectric, an atom
--- is itself a field carrier with its own mode structure. It imposes
boundary conditions: the tangential components of the electric field must
vanish at every conducting surface; the normal components of flux density
must be continuous across interfaces \EM. The total field is not the sum of
the partial fields --- it is the unique consistent field satisfying all
boundary conditions simultaneously.

Superposition is the approximation that arises when the boundary conditions
of individual objects are spatially far enough apart not to influence each
other. As soon as objects enter near-field range, this approximation breaks
down --- not gradually but structurally.

\subsection*{Near field: where boundary conditions collide}

The near-field zone of an antenna extends to roughly $r < 2D^2/\lambda$
(Fraunhofer boundary), where $D$ is the antenna size \EM. In this zone
the field cannot be described as a plane wave; it contains reactive
components (stored field energy oscillating back and forth without being
radiated). Bringing a second antenna into this zone causes the following:

\begin{itemize}[nosep]
\item The input impedance of both antennas changes --- the resonance
  frequency shifts, the radiation resistance changes. This is not field
  addition but mutual deformation of the mode structures.
\item The radiation pattern of each antenna changes because the boundary
  condition of the other deforms the field distribution in the shared
  near space.
\item Energy can be transferred directly between near fields --- without
  radiation into the far field. This is the principle of inductive coupling
  (transformer, Qi charging), capacitive coupling and resonant near-field
  coupling (Kurs et al.\ 2007 \cite{Kurs2007}, wireless transfer over metres).
\end{itemize}

\subsection*{Mutual coupling in antenna arrays}

In a phased array with element spacing $d < \lambda$ the mutual coupling
is not negligible \cite{Mailloux2017,Pozar2012}. The input impedance of the
$n$-th element depends on all others:
\[
  Z_{nn}^{\mathrm{active}} = Z_{nn}^{\mathrm{isolated}}
  + \sum_{m \neq n} Z_{nm}\,\frac{I_m}{I_n},
\]
where $Z_{nm}$ is the mutual impedance between elements $n$ and $m$. The
overall behaviour of the array is only describable by solving the complete
system of equations --- not by addition of single-element fields. In
practice mutual coupling is compensated by optimising element geometry and
feed phases; it is a design problem that industry solves with full
electromagnetic simulators (MoM, FEM, FDTD).

\subsection*{Cavity resonator: the total field as an eigenmode}

A metal cavity allows only discrete resonance frequencies determined by
its geometry \cite{Pozar2012,Griffiths1999}. For a rectangular cavity
of dimensions $a \times b \times c$:
\[
  f_{mnp} = \frac{c}{2}\sqrt{\left(\frac{m}{a}\right)^2
    + \left(\frac{n}{b}\right)^2 + \left(\frac{p}{c}\right)^2}.
\]
Placing an object inside the cavity --- a dielectric block, a conductor,
biological tissue --- shifts all resonance frequencies. The entire mode
spectrum of the cavity changes, not only the modes that directly touch the
object. The total field is the eigenmode of the composite system (cavity
+ object), not the sum of the empty-cavity eigenmode and a scattered field.
This shift is measurable and exploited in microwave dielectric spectroscopy
(cavity perturbation method).

\subsection*{Faraday cage: boundary conditions extinguish a field}

The Faraday cage is the most extreme case: a closed conducting shell
enforces vanishing tangential electric field on its inner surface. The
external field continues unchanged; the internal field is forced to zero
\EM. The total field is the unique solution of the Maxwell equations with
this boundary condition --- not superposition but cancellation through
enforced compatibility.

Incomplete shielding (gaps, finite conductivity) shows the principle in
attenuated form: fields penetrate but only with mode geometries permitted
by the gap geometry. A slot of width $w$ allows fields with $\lambda > 2w$
to pass through strongly attenuated (cut-off condition of the slot
waveguide). This is again mode compatibility: only fields whose wavelength
is compatible with the slot geometry couple through.

\subsection*{Plasmons: field and matter inseparable}

At the interface between metal and dielectric, surface plasmon polaritons
(SPP) arise --- coupled modes of electromagnetic field and collective
electron oscillation \cite{Maier2007}. The dispersion relation is:
\[
  k_\mathrm{SPP} = \frac{\omega}{c}\sqrt{\frac{\varepsilon_m\,\varepsilon_d}
    {\varepsilon_m + \varepsilon_d}},
\]
where $\varepsilon_m$ is the (complex) permittivity of the metal and
$\varepsilon_d$ that of the dielectric. The SPP is neither pure
electromagnetic field nor pure charge carrier oscillation --- it is a
joint mode of both carriers. Superposition would be meaningless: the
field exists only in coupling with matter.

This is the sharpest available evidence that ``multiple fields in space''
are not independent objects. The mode structure of the total system is
the physical reality; decomposition into partial fields is a computational
method that works under certain conditions.

\subsection*{FFGFT interpretation: one field, many projections}

In FFGFT there is only one field on the carrier $T^4/\mathbb{Z}_3$
\SetzM. What appears as ``multiple fields'' are projections of different
mode groups onto the observation space. The apparent superposition is a
representational question.

When two objects with different mode structures interact, the result is not
the sum of their fields but the joint eigenmode of the composite carrier
--- determined by the compatibility of the winding numbers of both carriers.
This is the same principle that appears in antenna engineering as mutual
coupling, in cavity physics as cavity perturbation and in nanophotonics
as plasmon coupling \SM.

The complexity of field interaction is therefore not a complication that
sufficient computational effort can remove --- it is the direct consequence
of every carrier imposing boundary conditions on the total field. The more
carriers in space, the more simultaneously required compatibility conditions,
the more complex the total field. This is not superposition --- this is
geometry.

% ============================================================
\section{Particles as field nodes: near and far field at localised field structures}
% ============================================================

In classical physics a particle is an object that \emph{has} a field ---
the charge generates the Coulomb field, the mass generates the gravitational
field. In FFGFT this description is inverted: a particle is a stable,
localised field configuration --- a node in the global field on the carrier
$T^4/\mathbb{Z}_3$ \SetzM. Mass is not an intrinsic property of the node
but follows from the fundamental equation $\tilde{T}\cdot m=1$: the time
scale of the field configuration determines the mass (Doc.~365 Layer~1,
Doc.~124) \KM.

This inversion has direct consequences for the structure of the field around
a particle --- and it makes the connection to antenna engineering visible.

\subsection*{Stationary particle: pure near field}

A stationary point charge $e$ generates the Coulomb field \EM:
\[
  \mathbf{E}(r) = \frac{e}{4\pi\varepsilon_0}\,\frac{\hat{r}}{r^2},
  \qquad
  \mathbf{B} = 0.
\]
In FFGFT units ($e=1$, $4\pi\varepsilon_0=1$): $V(r) = -1/r$
\cite{PascherDok365}. This field falls off as $1/r^2$ (field strength) or
$1/r$ (potential), carries no energy into the far field and is fully
reactive --- identical to the near field of a stationary antenna. The
Poynting vector $\mathbf{S} = \mathbf{E}\times\mathbf{H}$ is zero for a
stationary charge: no energy flux, only stored field energy in the node
structure.

The Coulomb field is the near field of the particle. It deforms the space
around the node, imposing boundary conditions on all other fields in its
vicinity --- exactly as a passive antenna deforms the field in its
surroundings without transmitting.

\subsection*{Accelerated particle: near field and radiated far field}

As soon as the charge is accelerated, a far field arises --- a detaching
electromagnetic wave. The Larmor formula gives the radiated power
\cite{Jackson1999,Griffiths1999}:
\[
  P = \frac{e^2 a^2}{6\pi\varepsilon_0 c^3}
    = \frac{2}{3}\,\frac{e^2 a^2}{c^3}
  \quad \text{(in Gaussian units)},
\]
where $a$ is the acceleration. In FFGFT units ($c=e=1$):
$P = \frac{2}{3}a^2$.

The far field falls off as $1/r$ (not $1/r^2$), permanently transports
energy away and is propagating --- identical to the far field of a driven
antenna. The boundary between near and far field lies at $r \sim \lambda/2\pi$
(transition zone), where $\lambda = c/f$ and $f$ is the characteristic
frequency of the acceleration.

This is the antenna principle at the particle level:
\begin{itemize}[nosep]
  \item Stationary charge = undriven antenna: near field only, no radiation.
  \item Uniformly moving charge = antenna on direct current: no alternating
    current, no radiation (Cherenkov radiation is a special case for $v > c/n$).
  \item Accelerated charge = driven antenna: near field plus radiated far
    field, power $\propto a^2$.
\end{itemize}

\subsection*{Atomic transition: quantised far field}

In the atom the permitted node structures are discrete through the quantum
numbers $(n, l, m)$ --- these are the permitted modes of the Coulomb
potential, i.e.\ the eigenmodes of the atomic carrier. A transition between
two states is a change in node structure; a photon is emitted --- the
radiated far field carries away a discrete energy $\Delta E = \hbar\omega$.
The selection rules \cite{CohenTannoudji1977,Griffiths1995}:
\[
  \Delta l = \pm 1, \qquad \Delta m = 0,\pm 1
\]
are compatibility conditions: only transitions in which the change of mode
geometry is compatible with the geometry of a dipole photon are permitted.
Quadrupole and higher transitions ($\Delta l = \pm 2$) are possible but
suppressed by a factor $\alpha^2 \approx 10^{-4}$ --- because the mode
compatibility is poorer.

In FFGFT language: the winding numbers of the initial and final state must
differ by exactly the winding numbers of a photon. The selection rules are
not postulated but follow from the geometry of the carrier \SM.

\subsection*{The ``size'' of a particle: extent of the node}

A particle as field node has no sharp geometric boundary. The ``size'' is
the spatial extent of the node structure, i.e.\ the region where the near
field dominates over the background field. For the electron these are the
classical electron radius \EM:
\[
  r_e = \frac{e^2}{4\pi\varepsilon_0 m_e c^2} \approx 2.82\,\mathrm{fm},
\]
the Compton wavelength $\lambda_C = \hbar/m_e c \approx 386\,\mathrm{fm}$
and the Bohr radius $a_0 \approx 52\,917\,\mathrm{fm}$. These three lengths
are not arbitrary but linked through $\alpha$ \cite{Jackson1999,PascherDok365}:
\[
  r_e = \alpha\,\lambda_C / 2\pi = \alpha^2 a_0.
\]
In FFGFT units ($\alpha=1$) these three lengths collapse to one --- the
hierarchy arises only through the SI embedding via $\xi$ (Doc.~365
Layer~3) \KM.

The node is therefore not point-like --- it has a structure extending over
several length scales. What appears as a ``point'' is the projection of
this structure onto a measurement coarse enough not to resolve the inner
structure.

\subsection*{Comparison table: antenna and particle}

\begin{center}
\begin{tabular}{@{}lll@{}}
\toprule
\textbf{Property} & \textbf{Antenna} & \textbf{Particle (node)} \\
\midrule
Near field & reactive, $E\propto 1/r^2$--$1/r^3$ & Coulomb field $E\propto 1/r^2$ \\
Far field & radiated wave, $E\propto 1/r$ & photon under acceleration \\
Radiation & only on current change & only under acceleration/transition \\
Mode structure & antenna geometry & quantum numbers $(n,l,m)$ \\
Selection rule & impedance compatibility & $\Delta l=\pm1$, $\Delta m=0,\pm1$ \\
``Size'' & physical dimension & node extent $r_e,\lambda_C,a_0$ \\
Stationary & no far field & no far field (Coulomb only) \\
Passive in field & scattering, shadowing & polarisability, cross-section \\
\bottomrule
\end{tabular}
\end{center}

\subsection*{Cross-section as a measure of node coupling}

A particle in an external field scatters the incoming wave --- it deforms
the field in its vicinity, exactly like a passive conductor. The Thomson
cross-section of the electron \cite{Jackson1999}:
\[
  \sigma_T = \frac{8\pi}{3}r_e^2 \approx 6.65\times 10^{-29}\,\mathrm{m}^2
\]
is the effective area from which the electron draws energy from the
incoming field and re-emits it. This is directly analogous to the effective
aperture $A_\mathrm{eff}$ of a receiving antenna (\S\,8):
\[
  A_\mathrm{eff} = \frac{\lambda^2}{4\pi}G
  \quad \longleftrightarrow \quad
  \sigma_T = \frac{8\pi}{3}r_e^2.
\]
In both cases the effective area is not the geometric size of the object
but the result of mode compatibility between incoming field and node
structure. For photons with $\hbar\omega \gg m_e c^2$, Compton scattering
applies --- the node structure of the electron can no longer be described
as point-like and the cross-section falls with increasing energy.

\subsection*{FFGFT interpretation: node as local mode concentration}

In FFGFT a particle is a local concentration of modes on the carrier
$T^4/\mathbb{Z}_3$ --- a location where specific winding numbers
$(n_\theta, n_\varphi)$ are occupied and the field locally assumes a
stable, non-trivial configuration \SetzM. The near field of the particle
is the immediate vicinity of this node in the carrier, projected onto the
observation space; the far field is the radiation arising when the node
changes its mode occupation.

The boundary between near and far field is therefore not a sharp line but
the transition between the region where the node structure dominates the
field and the region where the field propagates as a free wave. This
transition lies at $r\sim\lambda_C$ for relativistic descriptions and at
$r\sim a_0$ for atomic transitions --- both are structural length scales
of the node, not free parameters \SM.

% ============================================================
\section{Scale-invariance of the coupling condition: atoms, elements and planetary orbits}
% ============================================================

The principle from \S\,8--10 --- geometry of the carrier determines
permitted coupling angles, coupling only under mode compatibility --- is not
restricted to antennas or individual particles. It holds on all physical
scales. The scale changes; the structure of the condition remains identical.
Doc.~242 (scale ladder) elaborated this connection at the cosmological
level; here it is developed starting from antenna engineering.

\subsection*{Atoms: discrete modes in the Coulomb potential}

The Schr\"odinger equation in the Coulomb potential has discrete eigenmodes,
classified by quantum numbers $(n, l, m)$ \cite{Griffiths1995,CohenTannoudji1977}:
\[
  E_n = -\frac{m_e e^4}{2\hbar^2 n^2} = -\frac{13.6\,\mathrm{eV}}{n^2}.
\]
This is nothing other than the resonance condition of a carrier: the carrier
is the Coulomb potential (spherically symmetric, $V\propto -1/r$); the
permitted modes are the solutions with integer $n$. The atom is a cavity
resonator with spherical geometry --- the electron node can only assume
certain standing wave configurations, just as the TE$_{mn}$ modes in a
metal cavity (\S\,9).

The selection rules $\Delta l = \pm 1$, $\Delta m = 0,\pm 1$ (\S\,10) are
the impedance-matching conditions of this resonator: a dipole photon ($l=1$)
couples only to transitions in which the change of node geometry is
compatible with a dipole mode. A quadrupole photon ($l=2$) couples with
$\Delta l = \pm 2$ but more weakly by $\alpha^2$ --- because the mode
overlap efficiency is smaller, identical to the aperture efficiency
$\eta_A < 1$ of a horn antenna.

\subsection*{Periodic table: mode classification of a many-body system}

The shell structure of the periodic table --- $1s^2\,2s^2\,2p^6\,3s^2\ldots$
--- is a classification by mode geometries of the effective central
potential \EM:
\[
  \text{Shell } n: \quad 2n^2 \text{ places}
  \quad (s, p, d, f \text{ with spin}).
\]
Each sub-shell ($s, p, d, f$) has a different orbital angular momentum $l$
and thus a different spatial mode geometry. Chemical bonding arises through
coupling of the boundary modes of two atoms: when the outermost occupied
orbitals (HOMO/LUMO) have compatible geometry --- sufficient spatial
overlap, compatible symmetry --- a covalent bond forms. The overlap integral
\[
  S_{AB} = \int \psi_A^*(\mathbf{r})\,\psi_B(\mathbf{r})\,d^3r
\]
is the direct quantum-mechanical analogue of the polarisation factor
$\cos^2\psi$ in the Friis equation (\S\,8): $S_{AB}=0$ means no coupling,
regardless of the energy of the atoms.

The noble gases (closed shells, $S_{AB}\approx 0$ for all bonding partners)
are the quantum-mechanical analogue of orthogonal polarisation: no energy
transfer because no mode compatibility.

\subsection*{Crystal lattice: the Bragg condition as compatibility condition}

In a crystal the atoms form a periodic lattice of lattice constant $a$.
An electromagnetic wave (X-ray, neutrons, electrons) is coherently reflected
only when the Bragg condition is met \cite{Kittel2005,Ashcroft1976}:
\[
  2d\sin\theta = n\lambda, \quad n \in \mathbb{Z},
\]
where $d$ is the lattice plane spacing and $\theta$ the glancing angle.
This is an integer phase condition: the path-length difference between
partial waves reflected at successive lattice planes must be an integer
multiple of the wavelength. Only then do amplitudes add constructively ---
identical to the phased-array condition (\S\,8):
\[
  \phi_n = -n\,kd\cos\theta_0.
\]
The crystal lattice is a natural phased array with fixed element spacing
$d$. Bragg reflection selects exactly the angles of incidence at which mode
compatibility holds; all other angles give destructive interference. In
FFGFT the reciprocal lattice vectors $\mathbf{G}$ correspond to the
permitted winding vectors of the carrier; Bragg reflection is the
macroscopic manifestation of the quantisation condition on the torus \SM.

\subsection*{Molecular vibrations: infrared activity as a selection rule}

A molecular vibrational mode is infrared-active (absorbs IR radiation) only
when it produces a change in the electric dipole moment
\cite{Atkins2010,Wilson1955}:
\[
  \left(\frac{\partial \vec{\mu}}{\partial Q}\right)_{Q=0} \neq 0,
\]
where $Q$ is the normal coordinate of the vibration. This is the
impedance-matching condition between photon (dipole radiation) and
vibrational mode: only when the vibration has a dipole geometry does it
couple to a dipole photon. Symmetric stretching vibrations (as in N$_2$
or O$_2$) do not change the dipole moment --- they are IR-inactive but
Raman-active (coupling to quadrupole photons). This explains why N$_2$
and O$_2$ are not greenhouse gases: their vibrational modes are mode-incompatible
with IR dipole radiation.

\subsection*{Planetary orbits: integer resonances in the gravitational field}

The gravitational field of the sun is the carrier of the planetary system.
Its modes are the Kepler orbits, classified by the orbital parameters
$(a, e, i)$. When two planets are near an integer resonance $p:q$
($p, q \in \mathbb{Z}$), small gravitational perturbations accumulate
coherently over many revolutions (Doc.~242) \cite{PascherDok242,Murray1999}.
The Laplace resonance of the Galilean moons is the most famous example
\cite{Murray1999,Peale1979}:
\[
  \frac{T_\mathrm{Io}}{T_\mathrm{Europa}} \approx \frac{1}{2},
  \qquad
  \frac{T_\mathrm{Europa}}{T_\mathrm{Ganymede}} \approx \frac{1}{2},
  \qquad
  n_\mathrm{Io} - 3n_\mathrm{Europa} + 2n_\mathrm{Ganymede} = 0
\]
(exact relation of the mean motions $n = 2\pi/T$). This relation is not
a coincidence: it is a stable configuration because at exact resonance the
mutual perturbations cause no secular energy input, only periodic. The
system has an attractor at the integer resonance --- identical to the
attractor of an Arnold tongue in nonlinear dynamics (Doc.~328
\S\,coupling regimes) \KM.

Further examples:
\begin{itemize}[nosep]
  \item Jupiter/Saturn: $5:2$ resonance (``great inequality'',
    known since Laplace 1785) \QM.
  \item Neptune/Pluto: $3:2$ resonance (Plutinos) \QM.
  \item Kirkwood gaps in the asteroid belt: gaps at resonances $1:3$,
    $1:4$ with Jupiter (destructive resonance, not stabilising) \QM.
  \item Saturn's rings: gaps (Cassini Division) at $2:1$ resonance
    with Mimas \QM.
\end{itemize}

The coupling condition is the same as in antenna engineering: integer phase
compatibility. The carrier differs (gravitational rather than electromagnetic
field); the structure of the condition is identical.

\subsection*{Summary table: coupling condition on all scales}

\begin{center}
\begin{tabular}{@{}lllll@{}}
\toprule
\textbf{Scale} & \textbf{Carrier} & \textbf{Coupling condition}
  & \textbf{Example} & \textbf{Analogue} \\
\midrule
Antenna & radiator geometry & $\cos^2\psi$, impedance
  & dipole/Yagi & Friis equation \\
Atom & Coulomb potential & $\Delta l=\pm1$
  & H lines & selection rule \\
Molecule & eff.\ potential & $\partial\mu/\partial Q \neq 0$
  & IR activity & dipole coupling \\
Crystal & lattice periodicity & $2d\sin\theta=n\lambda$
  & Bragg reflection & phased array \\
Periodic table & shell structure & overlap $S_{AB}$
  & covalence & aperture overlap \\
Mol.\ bond & orbital geometry & symmetry compatibility
  & HOMO/LUMO & mode overlap \\
Planets & gravitational field & $p\,n_1 = q\,n_2$
  & Laplace resonance & Arnold tongue \\
\bottomrule
\end{tabular}
\end{center}

\subsection*{The common principle}

All rows of the table describe the same thing: a system has a carrier with
a mode spectrum; coupling to another system is only possible where the mode
spectra are compatible; incompatible modes do not couple, regardless of the
strength of the fields involved.

The technical maturity of antenna engineering --- decades of measurements,
standardised procedures, commercial simulators --- makes it the most
tangible realisation of this principle. Atomic mechanics, solid-state
physics and celestial mechanics are the same principle on different carriers.
In FFGFT the common carrier is $T^4/\mathbb{Z}_3$; the coupling conditions
on all scales are projections of winding-number compatibility onto the
respective observation space \SM.

% ============================================================
\section{Duality: field and particle, energy and current}
% ============================================================

FFGFT treats time--mass duality as a fundamental principle (fundamental
equation $\tilde{T}\cdot m=1$, Doc.~365) \SetzM. The same duality appears
in electromagnetic field theory:

\begin{center}
\begin{tabular}{@{}lll@{}}
\toprule
\textbf{Field view} & & \textbf{Current view} \\
\midrule
Poynting vector $\mathbf{S} = \mathbf{E}\times\mathbf{H}$ & $\longleftrightarrow$ & $P = V \cdot I$ \\
Field energy in dielectric & $\longleftrightarrow$ & $I^2 R$ Joule loss \\
Evanescent wave outside & $\longleftrightarrow$ & reactive power \\
Winding number (topological) & $\longleftrightarrow$ & mass (dynamical) \\
\bottomrule
\end{tabular}
\end{center}

In none of these pairs is one side ``more fundamental'' than the other ---
they describe the same field process from different projections. This is
the precise analogue of time--mass duality in FFGFT: $\tilde{T}$ and $m$
are reciprocal representations of the same object, not two different
quantities \KM.

\subsection*{Reactive power and evanescent fields}

On the G-line energy pulses back and forth in the periphery of the wire
field --- energy flux alternates direction in step with the field. This is
reactive power in AC circuit theory, or --- stated optically --- the
evanescent field on the shadow side of a totally reflecting interface \EM.
In FFGFT this behaviour corresponds to the non-propagating part of the
Fourier modes on $T^4/\mathbb{Z}_3$: modes below the spectral gap
(Doc.~314 \S\,H) carry no net energy but structure the spectrum \SM.

% ============================================================
\section{Exact agreement: Poynting integral and power formula}
% ============================================================

It is instructive to understand the exactness of this identity, because it
is regularly presented in courses as an ``approximate agreement'' or a
``plausibility argument''. Poynting's theorem states \EM:
\[
  \frac{\partial u}{\partial t} + \nabla \cdot \mathbf{S} = -\mathbf{j}\cdot\mathbf{E},
\]
where $u = \frac{1}{2}(\varepsilon_0 E^2 + \mu_0^{-1}B^2)$ is the
electromagnetic energy density. Integration over the volume of the cable
with the corresponding boundary conditions yields exactly $P = V \cdot I$
--- the only prerequisite is the validity of the Maxwell equations.

In FFGFT units (without $\varepsilon_0, \mu_0, e$) the identity reduces to:
\[
  P = E \cdot m,
\]
which follows immediately from $V = E$, $I = m$ and the fundamental
equation $\tilde{T}\cdot m=1$ \KM. The exactness is structural, not a
numerical coincidence.

% ============================================================
\section{Summary}
% ============================================================

\begin{enumerate}[nosep]
\item \textbf{Primary localisation.} Energy flows in the fields
  ($\mathbf{S} = \mathbf{E}\times\mathbf{H}$), not in the conductor.
  Voltage and current are macroscopic projections of the same field
  process \EM.

\item \textbf{Geometric objects.} Field lines on the carrier
  $T^4/\mathbb{Z}_3$ are topological objects with winding numbers; they
  classify modes and determine mass scales via $\tilde{T}\cdot m=1$ \BM.

\item \textbf{FFGFT units.} With $c=\hbar=\alpha=1$: $I=m$, $V=E$,
  $P=E\cdot m$ --- exactly, not approximately. The Maxwell equations
  carry no prefactors \KM.

\item \textbf{G-line as illustration.} The Goubau experiment demonstrates
  macroscopically that energy transport without an explicit return conductor
  is possible because the transport occurs in the field. The cone transition
  corresponds to impedance matching; field lines locally optimise energy
  flux \EM.

\item \textbf{Duality.} Field view and current view are reciprocal
  projections --- no pair is more fundamental. This mirrors the time--mass
  duality $\tilde{T}\cdot m=1$ at the level of classical fields \SM.
\end{enumerate}

\section*{References}
\addcontentsline{toc}{section}{References}

\begin{thebibliography}{99}

% --- Electromagnetism and antenna engineering ---

\bibitem{Goubau1950}
G.~Goubau,
\textit{Surface waves and their application to transmission lines},
J.~Appl.~Phys.\ \textbf{21} (1950) 1119--1128.
Foundational work on single-wire wave guidance; cone transition,
field configuration, loss measurement of 6~dB/mile.

\bibitem{Harms1907}
F.~Harms,
\textit{Elektromagnetische Wellen an einem Draht mit isolierender
zylindrischer H\"ulle},
Ann.~Phys.\ \textbf{23} (1907) 44--60.
First analysis of the dielectric as a field-guiding layer; precursor
of the Goubau line.

\bibitem{Stratton1941}
J.~A.~Stratton,
\textit{Electromagnetic Theory},
McGraw-Hill, New York, 1941.
Standard reference for the Poynting theorem (Ch.~2), boundary conditions,
field energy.

\bibitem{Griffiths1999}
D.~J.~Griffiths,
\textit{Introduction to Electrodynamics},
3rd ed., Prentice Hall, 1999.
Poynting vector, Larmor formula, radiation resistance, reciprocity;
accessible derivation of all standard results used in \S\,2--5.

\bibitem{Jackson1999}
J.~D.~Jackson,
\textit{Classical Electrodynamics},
3rd ed., Wiley, 1999.
Larmor formula (Ch.~14), Thomson cross-section (Ch.~14.8), classical
electron radius, Compton wavelength; canonical reference.

\bibitem{Balanis2016}
C.~A.~Balanis,
\textit{Antenna Theory: Analysis and Design},
4th ed., Wiley, 2016.
Standard text: radiation patterns, gain, radiation resistance, short dipole,
half-wave dipole, monopole, Yagi--Uda, horn, parabolic, patch, phased array,
Friis equation, reciprocity theorem.

\bibitem{Friis1946}
H.~T.~Friis,
\textit{A note on a simple transmission formula},
Proc.~IRE \textbf{34} (1946) 254--256.
Original paper on the link budget equation;
``\textit{The received power is [\ldots] proportional to the product of the
transmitting and receiving gains.}''

\bibitem{Yagi1928}
H.~Yagi,
\textit{Beam transmission of ultra short waves},
Proc.~IRE \textbf{16} (1928) 715--741.
Original Yagi--Uda paper; reflector and director principle;
gain through phase interference of passive elements.

\bibitem{Pozar2012}
D.~M.~Pozar,
\textit{Microwave Engineering},
4th ed., Wiley, 2012.
Cavity resonator (Ch.~6), patch antenna (Ch.~14), horn antenna,
waveguide modes, impedance matching, Smith chart, cavity perturbation.

\bibitem{Mailloux2017}
R.~J.~Mailloux,
\textit{Phased Array Antenna Handbook},
3rd ed., Artech House, 2017.
Array factor, beam steering, mutual coupling, active impedance matrix,
grating-lobe condition; quantitative treatment of all array phenomena.

\bibitem{Maier2007}
S.~A.~Maier,
\textit{Plasmonics: Fundamentals and Applications},
Springer, 2007.
Surface plasmon polaritons, dispersion relation (Ch.~2), coupling condition,
near-field enhancement; standard text of nanophotonics.

\bibitem{Kurs2007}
A.~Kurs, A.~Karalis, R.~Moffatt, J.~D.~Joannopoulos,
P.~Fisher, M.~Solja\v{c}i\'{c},
\textit{Wireless power transfer via strongly coupled magnetic resonances},
Science \textbf{317} (2007) 83--86.
Resonant near-field coupling over 2~m with 40\,\% efficiency;
``\textit{The transfer is mediated by the overlap of the near fields.}''
Direct evidence for energy transfer without radiation through mode
compatibility.

% --- Quantum mechanics and atomic physics ---

\bibitem{Griffiths1995}
D.~J.~Griffiths,
\textit{Introduction to Quantum Mechanics},
Prentice Hall, 1995.
Hydrogen atom, quantum numbers, selection rules $\Delta l=\pm1$,
$\Delta m=0,\pm1$ (Ch.~9); overlap integral and bonding theory (Ch.~7).

\bibitem{CohenTannoudji1977}
C.~Cohen-Tannoudji, B.~Diu, F.~Lalo\"e,
\textit{Quantum Mechanics}, 2~vols.,
Wiley, 1977.
Selection rules (Ch.~XIII), dipole approximation, quadrupole suppression
$\sim\alpha^2$; canonical reference for atomic transitions.

% --- Solid-state physics ---

\bibitem{Kittel2005}
C.~Kittel,
\textit{Introduction to Solid State Physics},
8th ed., Wiley, 2005.
Bragg condition (Ch.~2), reciprocal lattice, Brillouin zone, lattice
dynamics; standard solid-state text.

\bibitem{Ashcroft1976}
N.~W.~Ashcroft, N.~D.~Mermin,
\textit{Solid State Physics},
Holt, Rinehart and Winston, 1976.
Bragg scattering, phased-array analogy of reciprocal lattice (Ch.~6),
electron bands; advanced reference.

% --- Molecular physics ---

\bibitem{Atkins2010}
P.~Atkins, J.~de~Paula,
\textit{Physical Chemistry},
9th ed., Oxford University Press, 2010.
IR activity condition $\partial\vec{\mu}/\partial Q\neq0$ (Ch.~13),
Raman activity, symmetry selection rules; standard physical chemistry text.

\bibitem{Wilson1955}
E.~B.~Wilson, J.~C.~Decius, P.~C.~Cross,
\textit{Molecular Vibrations},
McGraw-Hill, 1955.
Normal coordinates, selection rules, group theory of vibrational modes;
classic original reference.

% --- Celestial mechanics ---

\bibitem{Murray1999}
C.~D.~Murray, S.~F.~Dermott,
\textit{Solar System Dynamics},
Cambridge University Press, 1999.
Orbital resonances (Ch.~8), Laplace resonance Io/Europa/Ganymede,
Kirkwood gaps, Cassini Division, Arnold tongues in celestial mechanics;
authoritative modern reference.

\bibitem{Peale1979}
S.~J.~Peale, P.~Cassen, R.~T.~Reynolds,
\textit{Melting of Io by tidal dissipation},
Science \textbf{203} (1979) 892--894.
Laplace resonance as the cause of tidal heating of Io;
``\textit{The forced eccentricity [\ldots] results from the Laplace relation
$n_1 - 3n_2 + 2n_3 = 0$.}'' Spectacular prediction before Voyager~1 arrival.

% --- FFGFT corpus (interlinks) ---

\bibitem{PascherDok365}
J.~Pascher,
\textit{Doc.~365: The Foundation of FFGFT in Four Layers ---
fundamental equation, electromagnetic sector, scaling parameter and
algebraic confirmations},
FFGFT corpus, Zenodo (2026).
Layer~1 ($\tilde{T}\cdot m=1$), Layer~2 (FFGFT units, $I=m$, $V=E$,
Maxwell without prefactors), Layer~3 ($\xi=4/30000$).

\bibitem{PascherDok314}
J.~Pascher,
\textit{Doc.~314: Lattice in Hilbert Space},
FFGFT corpus, Zenodo (2026).
Fourier modes on $T^4/\mathbb{Z}_3$, winding numbers as topological
invariants, shell count to norm~120, spectral gap.

\bibitem{PascherDok328}
J.~Pascher,
\textit{Doc.~328: Coupling regimes, resonance and synchronisation},
FFGFT corpus, Zenodo (2026).
Galois lattice condition, Arnold tongues, Adler model, $K_\mathrm{eff}=2\pi\xi$,
coupling threshold; direct connection to orbital resonances and mutual coupling.

\bibitem{PascherDok242}
J.~Pascher,
\textit{Doc.~242: Scale ladder --- interscale couplings in FFGFT},
FFGFT corpus, Zenodo (2026).
Planetary orbit argument, Laplace resonance as structural principle,
scale separation as historical convention, not physical necessity.

\bibitem{PascherDok124}
J.~Pascher,
\textit{Doc.~124: Charge as geometric projection},
FFGFT corpus, Zenodo (2026).
Charge as emergent projection of geometric patterns; particles as field
configurations, not intrinsic point properties.

\bibitem{PascherDok343}
J.~Pascher,
\textit{Doc.~343: The zeta function of the $T^4/\mathbb{Z}_3$ torus},
FFGFT corpus, Zenodo (2026).
Resolution floor $100\xi$, Galois lattice, Farey critical denominator
height; quantitative bound on mode distinguishability.

\end{thebibliography}
